\chapter*{序言：从精确方程到可信的近似}
\addcontentsline{toc}{chapter}{序言：从精确方程到可信的近似}

高级物理最容易给人一种错觉：章节越往后，公式只是越来越长。真正增加的并不是代数长度，而是同一个问题里同时出现了更多尺度。有限晶格与热力学极限、薄板厚度与面内波长、环境记忆时间与系统弛豫时间、电子运动与核运动，都是因为两个尺度相差很大，才允许我们用一个较简单的理论替代完整动力学。一个近似是否可信，取决于被省略的量究竟有多小，而不是它在教科书里有多常见。

以相变为例。任何有限系统的配分函数都是良性的有限和，真正的非解析性只能在自由度趋于无穷时出现。Lee--Yang 零点、图展开和转移算子看似是三种不同技巧，其实都在追踪有限系统怎样逼近一个可能不解析的极限。薄板理论也有同样的逻辑：三维弹性方程并没有消失，只是在厚度 \(h\) 远小于面内尺度 \(L\) 时，横向剪切和厚度伸缩退到更高阶，弯曲能才留下四阶的中面方程。

量子理论中的约化更加需要这种谨慎。对系统与环境的联合态取偏迹是精确操作；Born 近似、马尔可夫近似和久期近似却分别删去不同的信息。把它们统称为“弱耦合”会掩盖真正的物理条件。分子绝热近似同样不是一句“核比电子重”：电子本征空间随核构型移动，由此产生的导数耦合、Berry 联络和几何标量势必须先保留下来，能隙与核速度的比较才决定哪些通道可以舍去。

本书中的推导因而遵守一个简单习惯：先写出尚未近似的对象，说明它生活在哪个空间、满足什么边界条件或守恒律；再引入无量纲小参数；最后检查近似在极限之外怎样失效。这样处理时，统计物理、连续介质、开放系统、原子光学和分子动力学并不是互不相干的专题。它们都在练习同一件事——怎样从精确结构中取出可计算的有效理论，同时把被舍弃的信息和适用范围留在明处。

公式当然重要，但公式之间的句子同样重要。读者应当能够在每一次计算前说清要回答的问题，在计算后说清结果改变了什么。若一条等式只能被抄写而不能被解释，它还没有真正成为理论的一部分。
